% TOP_PROBLEM: {"id": 2305016, "title": "Research Problems in Function Theory — Problem 5.16", "category_id": 8, "category": {"id": 8, "name": "analysis", "display_name": "Analysis", "description": "Limits, continuity, calculus, and function theory.", "slug": "analysis", "order_index": 8, "created_at": "2026-07-31T15:26:25.671Z"}, "status": "open"}
% FIRST_POSTED: null
% ENTRY_KIND: statement_only
% Imported from: batch14.tex; UnsolvedMath ID 2305016
% Compile twice: lualatex 2305016.tex
\documentclass[11pt]{article}
\usepackage{fontspec}
\setmainfont{Latin Modern Roman}
\usepackage{amsmath,amssymb,mathtools,unicode-math}
\setmathfont{Latin Modern Math}
\usepackage{xurl,hyperref}
\hypersetup{hidelinks,pdftitle={UnsolvedMath Problem 2305016}}
\newcommand{\sref}[2]{\hyperlink{source:#1:#2}{[S#2]}}
\newcommand{\cataloguescope}[1]{\paragraph{Mathematical question.}#1}
\begin{document}
% UnsolvedMath ID 2305016; source catalogue code AMR-022-5016
\setcounter{section}{2305015}
\section{Research Problems in Function Theory — Problem 5.16}\label{2305016}
{\small\sffamily UnsolvedMath ID 2305016 · Analysis}
\subsection{Definitions and mathematical statement}

\paragraph{Shared notation.}$\mathbb N=\{1,2,\ldots\}$, and $\mathbb Z,\mathbb Q,\mathbb R,\mathbb C$ denote the integers, rationals, reals and complex numbers. Any other conventions are specified locally.


Write $\mathbb D=\{z\in\mathbb C:|z|<1\}$ and $\mathbb T=\partial\mathbb D$. All Taylor coefficients below are taken at $0$. For a plane domain $D$, write $\ell_D(r)$ for the angular Lebesgue measure of $\{\theta:re^{i\theta}\in D\}$. Its circular symmetrization $D^*$ contains the whole circle of radius $r$ when $D$ does; contains no points on that circle when $D$ misses it; otherwise it contains the arc $\{re^{i\theta}:|\theta|<\ell_D(r)/2\}$. Include $0$ in $D^*$ exactly when $0\in D$. Let $a_0>0$ belong to $D$, let $f:\mathbb D\to D$ be holomorphic with $f(0)=a_0$, and suppose $D^*$ admits a universal covering map $f^*:\mathbb D\to D^*$. Normalize the covering map by $f^*(0)=a_0$ and $(f^*)'(0)>0$. Write $f(z)=a_0+\sum_{n\ge1}a_nz^n$ and $f^*(z)=a_0+\sum_{n\ge1}a_n^*z^n$.Define, for $0<r<1$ and $\lambda>0$,
\[I_\lambda(r,h)=\left(\frac1{2\pi}\int_0^{2\pi}|h(re^{i\theta})|^\lambda\,d\theta\right)^{1/\lambda},\qquad
T(r,h)=\frac1{2\pi}\int_0^{2\pi}\log^+|h(re^{i\theta})|\,d\theta,\]
where $\log^+x=\max\{0,\log x\}$, with $\log^+0=0$. For holomorphic $h$, $T$ is its Nevanlinna characteristic.
\paragraph{Statement recorded in the supplied source.}
For arbitrary domains in this setup, without a simply-connectedness assumption, do
\[I_\lambda(r,f)\le I_\lambda(r,f^*)\quad\text{for every }\lambda>0,\qquad
T(r,f)\le T(r,f^*)\]
hold for all $0<r<1$? \sref{2305016}{2}
\subsection{Short English statement}
Does the universal cover of the circularly symmetrized image domain dominate every positive integral mean and the Nevanlinna characteristic?
\subsection{Sources}
\begin{description}
\item[\hypertarget{source:2305016:1}{S1}] ULAM.AI and the UnsolvedMath Contributors, UnsolvedMath: A Curated Collection of Open Mathematics Problems, record 2305016 (AMR-022-5016). CC BY 4.0. \url{https://huggingface.co/datasets/ulamai/UnsolvedMath}.
\item[\hypertarget{source:2305016:2}{S2}] W. K. Hayman and E. F. Lingham, Research Problems in Function Theory (2018), arXiv:1809.07200, Chapter 5, Problem 5.16, its update and the preceding shared notation. \url{https://arxiv.org/abs/1809.07200}.
\end{description}
\label{end:2305016}
\end{document}
